% Auto-generated from raw.txt
% Usage:
%   \vocabentry{word}{/ipa/ (pos)}{definition}{example}{note}

\vocabentry{Urbanization}
{/ˌɜːrbənəˈzeɪʃn/ (n)}
{The process of making an area more urban.}
{Rapid urbanization in developing countries often leads to a shortage of affordable housing.}
{Đây là một danh từ dùng để chỉ quá trình đô thị hóa, khi dân cư tập trung đông đúc vào các thành phố.}

\vocabentry{Infrastructure}
{/ˈɪnfrəstrʌktʃər/ (n)}
{The basic physical and organizational structures needed for the operation of a society.}
{The government is investing billions in urban infrastructure to reduce traffic congestion.}
{Đây là một danh từ dùng để chỉ cơ sở hạ tầng như cầu đường, hệ thống điện nước.}

\vocabentry{Metropolis}
{/məˈtrɑːpəlɪs/ (n)}
{The capital or chief city of a country or region.}
{Tokyo is a vast metropolis that blends traditional culture with futuristic technology.}
{Đây là một danh từ dùng để chỉ một thành phố lớn, trung tâm kinh tế và văn hóa quan trọng.}

\vocabentry{Congestion}
{/kənˈdʒestʃən/ (n)}
{The state of being extremely full or blocked with something, typically traffic.}
{Traffic congestion is a major problem that affects the quality of life in big cities.}
{Đây là một danh từ dùng để chỉ sự tắc nghẽn, đặc biệt là tắc nghẽn giao thông.}

\vocabentry{Cosmopolitan}
{/ˌkɑːzməˈpɑːlɪtən/ (adj)}
{Including people from many different countries and cultures.}
{London is a truly cosmopolitan city where you can hear hundreds of different languages.}
{Đây là một tính từ dùng để chỉ tính đa văn hóa hoặc mang tầm quốc tế của một thành phố.}

\vocabentry{Commute}
{/kəˈmjuːt/ (v)}
{To travel some distance between one's home and place of work on a regular basis.}
{Many people choose to live in the suburbs and commute to the city center for work.}
{Đây là một động từ dùng để chỉ việc đi làm hàng ngày từ nhà đến cơ quan và ngược lại.}

\vocabentry{Affordable}
{/əˈfɔːrdəbl/ (adj)}
{Inexpensive; reasonably priced.}
{The lack of affordable housing is forcing young professionals to move further away from the city.}
{Đây là một tính từ dùng để chỉ giá cả phải chăng, thường dùng trong cụm "affordable housing" - nhà ở giá rẻ.}

\vocabentry{Amenity}
{/əˈmenəti/ (n)}
{A desirable or useful feature or facility of a building or place.}
{The new residential complex offers a wide range of amenities, including a swimming pool and a gym.}
{Đây là một danh từ dùng để chỉ các tiện nghi công cộng như công viên, rạp phim, trung tâm thể thao.}

\vocabentry{Slum}
{/slʌm/ (n)}
{A squalid and overcrowded urban street or district inhabited by very poor people.}
{The government is working on a project to provide better sanitation for people living in the city's slums.}
{Đây là một danh từ dùng để chỉ khu ổ chuột, nơi điều kiện sống thấp kém.}

\vocabentry{Suburb}
{/ˈsʌbɜːrb/ (n)}
{An outlying district of a city, especially a residential one.}
{They decided to move to the suburbs to raise their children in a cleaner environment.}
{Đây là một danh từ dùng để chỉ khu vực ngoại ô, nơi yên tĩnh hơn trung tâm thành phố.}

\vocabentry{Gentrification}
{/ˌdʒentrɪfɪˈkeɪʃn/ (n)}
{The process of changing the character of a neighborhood through the influx of more affluent residents and businesses.}
{Gentrification has transformed this once-neglected district into a trendy area with expensive boutiques.}
{Đây là một danh từ dùng để chỉ sự chỉnh trang đô thị, thường khiến giá nhà tăng cao và đẩy người nghèo đi nơi khác.}

\vocabentry{Density}
{/ˈdensəti/ (n)}
{The degree of compactness of a substance; the number of people in a given area.}
{The high population density in Hong Kong makes land extremely valuable.}
{Đây là một danh từ dùng để chỉ mật độ, cụ thể là mật độ dân số.}

\vocabentry{Sprawl}
{/sprɔːl/ (n)}
{The disorganized and inefficient expansion of an urban area into the surrounding countryside.}
{Urban sprawl leads to the destruction of agricultural land and increased dependence on cars.}
{Đây là một danh từ dùng để chỉ sự mở rộng đô thị tràn lan, thiếu kiểm soát.}

\vocabentry{Inhabitant}
{/ɪnˈhæbɪtənt/ (n)}
{A person or animal that lives in or occupies a place.}
{The city's inhabitants are demanding better public transportation and more green spaces.}
{Đây là một danh từ dùng để chỉ cư dân sinh sống tại một địa điểm cụ thể.}

\vocabentry{Hectic}
{/ˈhektɪk/ (adj)}
{Full of incessant or frantic activity.}
{I need a holiday to get away from the hectic pace of city life for a while.}
{Đây là một tính từ dùng để chỉ nhịp sống bận rộn, hối hả và căng thẳng ở thành phố.}

\vocabentry{Modernization}
{/ˌmɑːrdənəˈzeɪʃn/ (n)}
{The process of adapting something to modern needs or habits.}
{The modernization of the old railway station has made traveling much more comfortable.}
{Đây là một danh từ dùng để chỉ sự hiện đại hóa cơ sở vật chất hoặc lối sống.}

\vocabentry{Skyscraper}
{/ˈskaɪskreɪpər/ (n)}
{A very tall continuous building having many floors.}
{The city skyline is dominated by impressive skyscrapers made of glass and steel.}
{Đây là một danh từ dùng để chỉ những tòa nhà chọc trời, biểu tượng của các đô thị lớn.}

\vocabentry{Pavement}
{/ˈpeɪvmənt/ (n)}
{A raised paved or asphalted path for pedestrians at the side of a road.}
{The city council is repairing the broken pavements to improve pedestrian safety.}
{Đây là một danh từ dùng để chỉ vỉa hè dành cho người đi bộ.}

\vocabentry{Residential}
{/ˌrezɪˈdenʃl/ (adj)}
{Designed for people to live in.}
{This is a quiet residential neighborhood with very little through traffic.}
{Đây là một tính từ dùng để chỉ các khu vực dùng để ở, thay vì dùng cho công nghiệp hay thương mại.}

\vocabentry{Sanitation}
{/ˌsænɪˈteɪʃn/ (n)}
{Conditions relating to public health, especially the provision of clean drinking water and adequate sewage disposal.}
{Poor sanitation in overcrowded cities can lead to the rapid spread of infectious diseases.}
{Đây là một danh từ dùng để chỉ hệ thống vệ sinh công cộng.}

\vocabentry{Commercial}
{/kəˈmɜːrʃl/ (adj)}
{Concerned with or engaged in commerce.}
{The downtown area is primarily a commercial district filled with office buildings and shops.}
{Đây là một tính từ dùng để chỉ các khu vực hoặc hoạt động mang tính thương mại, kinh doanh.}

\vocabentry{Overcrowding}
{/ˌoʊvərˈkraʊdɪŋ/ (n)}
{The presence of too many people in one place.}
{Overcrowding on public buses is a common complaint among city residents.}
{Đây là một danh từ dùng để chỉ tình trạng đông đúc quá tải.}

\vocabentry{Redevelopment}
{/ˌriːdɪˈveləpmənt/ (n)}
{The action or process of developing something again or differently.}
{The redevelopment of the old dockyards has created thousands of new jobs.}
{Đây là một danh từ dùng để chỉ việc quy hoạch hoặc xây dựng lại một khu vực đô thị.}

\vocabentry{Exodus}
{/ˈeksədəs/ (n)}
{A mass departure of people.}
{The rural-to-urban exodus has put immense pressure on city services.}
{Đây là một danh từ dùng để chỉ sự di cư của một nhóm người lớn, ví dụ từ nông thôn ra thành thị.}

\vocabentry{Slumdog}
{/ˈslʌmdɔːɡ/ (n)}
{A person who lives in a slum. (Đây là một danh từ (không trang trọng) dùng để chỉ người cư ngụ ở khu ổ chuột.).}
{The movie depicts the struggles and resilience of a slumdog in Mumbai.}
{}

\vocabentry{Underground}
{/ˌʌndərˈɡraʊnd/ (n)}
{An electric passenger railway in a city operating in tunnels below the ground.}
{Taking the underground is usually the fastest way to get across the city during rush hour.}
{Đây là một danh từ dùng để chỉ hệ thống tàu điện ngầm.}

\vocabentry{Skyline}
{/ˈskaɪlaɪn/ (n)}
{An outline of land and buildings defined against the sky.}
{New York has one of the most recognizable skylines in the world.}
{Đây là một danh từ dùng để chỉ đường chân trời, thường là hình dáng các tòa nhà khi nhìn từ xa.}

\vocabentry{Vibrancy}
{/ˈvaɪbrənsi/ (n)}
{The state of being full of energy and life.}
{I love the cultural vibrancy of New Orleans, especially during the jazz festival.}
{Đây là một danh từ dùng để chỉ sự náo nhiệt, tràn đầy sức sống của một thành phố.}

\vocabentry{Pollution}
{/pəˈluːʃn/ (n)}
{The presence in or introduction into the environment of a substance which has harmful or poisonous effects.}
{Air pollution in megacities has reached dangerous levels in recent years.}
{Đây là một danh từ dùng để chỉ sự ô nhiễm không khí, nước hoặc tiếng ồn ở đô thị.}

\vocabentry{Megacity}
{/ˈmeɡəsɪti/ (n)}
{A very large city, typically one with a population of over ten million people.}
{By 2050, more than half of the world's population will live in megacities.}
{Đây là một danh từ dùng để chỉ siêu đô thị với dân số cực lớn.}

\vocabentry{Conurbation}
{/ˌkɑːnɜːrˈbeɪʃn/ (n)}
{An extended urban area, typically consisting of several towns merging with the suburbs of one or more cities.}
{The area between Tokyo and Yokohama has become one continuous conurbation.}
{Đây là một danh từ dùng để chỉ vùng đô thị mở rộng, gồm nhiều thành phố nhỏ kết nối lại.}

\vocabentry{Bustling}
{/ˈbʌslɪŋ/ (adj)}
{Full of energetic and noisy activity.}
{The bustling markets of Marrakesh are a feast for the senses.}
{Đây là một tính từ dùng để chỉ sự nhộn nhịp, tấp nập của phố xá.}

\vocabentry{Pedestrianization}
{/pəˌdestriənaɪˈzeɪʃn/ (n)}
{The process of making a street or area for use by people walking only.}
{The pedestrianization of the city center has encouraged more people to shop and dine outdoors.}
{Đây là một danh từ dùng để chỉ việc chuyển một khu phố thành phố đi bộ.}

\vocabentry{Inner-city}
{/ˌɪnər ˈsɪti/ (n)}
{The area near the center of a large city, often associated with social and economic problems.}
{Many inner-city schools are struggling with a lack of funding and resources.}
{Đây là một danh từ dùng để chỉ khu vực nội thành, thường là nơi dân cư nghèo sinh sống.}

\vocabentry{Ghetto}
{/ˈɡetoʊ/ (n)}
{A part of a city, especially a slum area, occupied by a minority group or groups.}
{The city is working to integrate the residents of the ghetto into the wider community.}
{Đây là một danh từ dùng để chỉ khu vực tập trung các nhóm thiểu số, thường có điều kiện sống khó khăn.}

\vocabentry{Livable}
{/ˈlɪvəbl/ (adj)}
{(Of a city or environment) fit to live in; enjoyable.}
{Vienna is consistently ranked as one of the most livable cities in the world.}
{Đây là một tính từ dùng để chỉ tính đáng sống của một thành phố.}

\vocabentry{Downtown}
{/ˌdaʊnˈtaʊn/ (n)}
{The central part or main business and commercial area of a town or city.}
{Most of the major banks and corporations have their headquarters downtown.}
{Đây là một danh từ dùng để chỉ trung tâm thành phố, nơi tập trung các hoạt động kinh tế.}

\vocabentry{Fringe}
{/frɪndʒ/ (n)}
{The outer edge or margin of an area.}
{New housing developments are springing up on the rural-urban fringe.}
{Đây là một danh từ dùng để chỉ vùng rìa của đô thị.}

\vocabentry{Decentralization}
{/ˌdiːˌsentrələˈzeɪʃn/ (n)}
{The transfer of authority from central to local government or the movement of departments from a city center to the outskirts.}
{Decentralization can help reduce the pressure on overcrowded city centers.}
{Đây là một danh từ dùng để chỉ sự phi tập trung hóa, dời các hoạt động ra xa lõi thành phố.}

\vocabentry{Industrial}
{/ɪnˈdʌstriəl/ (adj)}
{Relating to or characterized by industry.}
{The old industrial zone is being converted into modern loft apartments.}
{Đây là một tính từ dùng để chỉ các khu công nghiệp, thường nằm ở xa trung tâm dân cư.}

\vocabentry{Zoning}
{/ˈzoʊnɪŋ/ (n)}
{The process of dividing land in a municipality into zones in which certain uses are permitted or prohibited.}
{Proper zoning laws are essential to prevent factories from being built in residential areas.}
{Đây là một danh từ dùng để chỉ việc phân vùng quy hoạch sử dụng đất.}

\vocabentry{Traffic-free}
{/ˈtræfɪk friː/ (adj)}
{(Of an area) restricted to pedestrians and cyclists; no cars allowed.}
{The town square is a traffic-free zone, making it very safe for children to play.}
{Đây là một tính từ dùng để chỉ khu vực không có xe cộ, chỉ dành cho người đi bộ.}

\vocabentry{Concrete jungle}
{/ˌkɑːŋkriːt ˈdʒʌŋɡl/ (n)}
{A modern city or urban area which has a large number of large buildings and is perceived as unpleasant.}
{After living in the concrete jungle for years, she decided to move to a farm.}
{Đây là một cụm danh từ mang nghĩa bóng chỉ thành phố toàn bê tông, ngột ngạt và thiếu mảng xanh.}

\vocabentry{High-rise}
{/ˈhaɪ raɪz/ (n)}
{A building with many stories.}
{Living in a high-rise offers great views but can feel isolating for some people.}
{Đây là một danh từ dùng để chỉ các tòa nhà cao tầng.}

\vocabentry{Multi-storey}
{/ˌmʌlti ˈstɔːri/ (adj)}
{(Of a building) having several stories.}
{A new multi-storey car park is being built to address the shortage of parking spaces.}
{Đây là một tính từ dùng để chỉ các công trình nhiều tầng, thường là bãi đậu xe.}

\vocabentry{Congestion charge}
{/kənˈdʒestʃən tʃɑːrdʒ/ (n)}
{An amount of money that people must pay to drive into a city center.}
{London introduced a congestion charge to discourage car use and reduce pollution.}
{Đây là một cụm danh từ dùng để chỉ phí ùn tắc, đánh vào xe cá nhân đi vào trung tâm.}

\vocabentry{Mixed-use}
{/ˌmɪkst ˈjuːs/ (adj)}
{Combining residential, commercial, and industrial uses in a single development.}
{Mixed-use developments promote walkability by keeping homes and workplaces close together.}
{Đây là một tính từ dùng để chỉ các khu quy hoạch phức hợp, vừa có nhà ở vừa có văn phòng và cửa hàng.}

\vocabentry{Connectivity}
{/ˌkɑːnekˈtɪvəti/ (n)}
{The state of being connected or interconnected.}
{Excellent connectivity between the airport and the city center is vital for business.}
{Đây là một danh từ dùng để chỉ khả năng kết nối về giao thông hoặc internet trong đô thị.}

\vocabentry{Gated community}
{/ˌɡeɪtɪd kəˈmjuːnəti/ (n)}
{A residential area with restricted access and monitored by security.}
{Gated communities offer high levels of security and privacy for their residents.}
{Đây là một cụm danh từ dùng để chỉ khu dân cư có cổng bảo vệ, thường dành cho người giàu.}

\vocabentry{Public transport}
{/ˌpʌblɪk ˈtrænspɔːrt/ (n)}
{Buses, trains, and other forms of transport that are available to the public.}
{Investing in reliable public transport is the best way to reduce traffic jams.}
{Đây là một cụm danh từ dùng để chỉ phương tiện giao thông công cộng.}

\vocabentry{Nightlife}
{/ˈnaɪtlaɪf/ (n)}
{Entertainment available at night in a town or city.}
{Berlin is famous for its diverse and exciting nightlife, with clubs that stay open until dawn.}
{Đây là một danh từ dùng để chỉ các hoạt động vui chơi giải trí về đêm.}

\vocabentry{Street vendor}
{/ˈstriːt vendər/ (n)}
{A person who sells food or other goods in the street.}
{Street vendors provide affordable and delicious meals for busy office workers.}
{Đây là một cụm danh từ dùng để chỉ người bán hàng rong.}

\vocabentry{Sanitary}
{/ˈsænɪteri/ (adj)}
{Relating to the conditions that affect hygiene and health.}
{Maintaining sanitary conditions in public markets is essential to prevent food poisoning.}
{Đây là một tính từ dùng để chỉ các điều kiện vệ sinh.}

\vocabentry{Pedestrian}
{/pəˈdestriən/ (n)}
{A person walking rather than traveling in a vehicle.}
{The new bridge was designed specifically for pedestrians and cyclists.}
{Đây là một danh từ dùng để chỉ người đi bộ.}

\vocabentry{Cycle lane}
{/ˈsaɪkl leɪn/ (n)}
{A part of a road marked off for use by people riding bicycles.}
{Adding more cycle lanes to city streets encourages people to use eco-friendly transport.}
{Đây là một cụm danh từ dùng để chỉ làn đường dành riêng cho xe đạp.}

\vocabentry{Green space}
{/ˈɡriːn speɪs/ (n)}
{An area of grass, trees, or other vegetation set apart for recreational purposes in an urban environment.}
{Green spaces are essential for the mental health and well-being of city dwellers.}
{Đây là một cụm danh từ dùng để chỉ mảng xanh đô thị như công viên.}

\vocabentry{Urban dweller}
{/ˌɜːrbən ˈdwelər/ (n)}
{A person who lives in a city or town.}
{Urban dwellers often suffer from higher levels of stress than those living in rural areas.}
{Đây là một cụm danh từ dùng để chỉ cư dân thành phố.}

\vocabentry{Accessibility}
{/əkˌsesəˈbɪləti/ (n)}
{The quality of being able to be reached or entered.}
{The accessibility of the new library for disabled people was a key design priority.}
{Đây là một danh từ dùng để chỉ khả năng tiếp cận các dịch vụ công cộng.}

\vocabentry{Social isolation}
{/ˌsoʊʃl ˌaɪsəˈleɪʃn/ (n)}
{A state of complete or near-complete lack of contact between an individual and society.}
{Modern city life can sometimes lead to social isolation, despite the crowds.}
{Đây là một cụm danh từ dùng để chỉ sự cô lập xã hội, thường gặp ở người già tại các đô thị lớn.}

\vocabentry{Sense of community}
{/sens əv kəˈmjuːnəti/ (n)}
{A feeling that members have of belonging.}
{Smaller neighborhoods often have a stronger sense of community than large apartment blocks.}
{Đây là một cụm danh từ dùng để chỉ cảm giác thuộc về một cộng đồng, sự gắn kết giữa hàng xóm.}

\vocabentry{Standard of living}
{/ˌstændərd əv ˈlɪvɪŋ/ (n)}
{The degree of wealth and material comfort available to a person or community.}
{Urbanization has significantly raised the standard of living for many millions of people.}
{Đây là một cụm danh từ dùng để chỉ mức sống.}

\vocabentry{Cost of living}
{/ˌkɔːst əv ˈlɪvɪŋ/ (n)}
{The amount of money needed to sustain a certain level of living.}
{The high cost of living in cities like San Francisco makes it difficult for service workers to survive.}
{Đây là một cụm danh từ dùng để chỉ chi phí sinh hoạt.}

\vocabentry{Employment opportunity}
{/ɪmˈplɔɪmənt ˌɑːpərˈtuːnəti/ (n)}
{A chance to have a job.}
{People move to cities primarily in search of better employment opportunities.}
{Đây là một cụm danh từ dùng để chỉ cơ hội việc làm.}

\vocabentry{Social mobility}
{/ˌsoʊʃl moʊˈbɪləti/ (n)}
{The movement of individuals or groups in social position over time.}
{Education is the most effective tool for promoting social mobility in urban areas.}
{Đây là một cụm danh từ dùng để chỉ sự dịch chuyển tầng lớp xã hội.}

\vocabentry{Urban planning}
{/ˌɜːrbən ˈplænɪŋ/ (n)}
{The technical and political process concerned with the development and design of land use.}
{Good urban planning can create cities that are both efficient and beautiful.}
{Đây là một cụm danh từ dùng để chỉ quy hoạch đô thị.}

\vocabentry{Noise pollution}
{/ˈnɔɪz pəˌluːʃn/ (n)}
{Harmful or annoying levels of noise.}
{Noise pollution from constant construction is a major source of stress for residents.}
{Đây là một cụm danh từ dùng để chỉ ô nhiễm tiếng ồn.}

\vocabentry{Light pollution}
{/ˈlaɪt pəˌluːʃn/ (n)}
{Brightening of the night sky caused by street lights and other man-made sources.}
{Light pollution in cities makes it impossible to see most of the stars at night.}
{Đây là một cụm danh từ dùng để chỉ ô nhiễm ánh sáng.}

\vocabentry{Heat island}
{/ˈhiːt ˌaɪlənd/ (n)}
{An urban area that is significantly warmer than its surrounding rural areas due to human activities.}
{Planting trees on rooftops can help cool down city heat islands.}
{Đây là một cụm danh từ dùng để chỉ hiện tượng đảo nhiệt đô thị.}

\vocabentry{Sanitary landfill}
{/ˌsænɪteri ˈlændfɪl/ (n)}
{A site for the disposal of waste materials.}
{The city is struggling to find a suitable location for a new sanitary landfill.}
{Đây là một cụm danh từ dùng để chỉ bãi rác hợp vệ sinh.}

\vocabentry{Waste disposal}
{/ˈweɪst dɪˌspoʊzl/ (n)}
{The act of getting rid of waste material.}
{Effective waste disposal systems are essential to maintain urban hygiene.}
{Đây là một cụm danh từ dùng để chỉ việc xử lý rác thải.}

\vocabentry{Recycling scheme}
{/ˌriːˈsaɪklɪŋ skiːm/ (n)}
{A system for collecting and processing materials that would otherwise be thrown away as trash.}
{The local council has introduced a mandatory recycling scheme for all households.}
{Đây là một cụm danh từ dùng để chỉ kế hoạch/hệ thống tái chế.}

\vocabentry{Water supply}
{/ˈwɔːtər səˌplaɪ/ (n)}
{The provision of water by public utilities.}
{A reliable water supply is a fundamental requirement for any urban population.}
{Đây là một cụm danh từ dùng để chỉ nguồn cung cấp nước.}

\vocabentry{Gridlock}
{/ˈɡrɪdlɑːk/ (n)}
{A traffic jam in which no vehicle can move in any direction.}
{The accident during rush hour caused total gridlock in the city center.}
{Đây là một danh từ dùng để chỉ tình trạng giao thông tê liệt hoàn toàn.}

\vocabentry{One-way street}
{/ˌwʌn ˈweɪ striːt/ (n)}
{A street where traffic is allowed to move in only one direction.}
{Converting more roads into one-way streets can help improve the flow of traffic.}
{Đây là một cụm danh từ dùng để chỉ đường một chiều.}

\vocabentry{Flyover}
{/ˈflaɪoʊvər/ (n)}
{A bridge that carries one road over another.}
{The new flyover has significantly reduced the travel time between the two districts.}
{Đây là một danh từ dùng để chỉ cầu vượt.}

\vocabentry{Cul-de-sac}
{/ˈkʌl də sæk/ (n)}
{A street or passage closed at one end.}
{They live in a quiet cul-de-sac, so there is very little noise from cars.}
{Đây là một danh từ dùng để chỉ đường cụt, thường ở các khu dân cư yên tĩnh.}

\vocabentry{Landmark}
{/ˈlændmɑːrk/ (n)}
{An object or feature of a landscape or town that is easily seen and recognized from a distance.}
{The Eiffel Tower is the most famous landmark in Paris.}
{Đây là một danh từ dùng để chỉ công trình biểu tượng hoặc địa danh dễ nhận biết.}

\vocabentry{District}
{/ˈdɪstrɪkt/ (n)}
{An area of a country or city, especially one characterized by a particular feature or activity.}
{The financial district is deserted on weekends when the offices are closed.}
{Đây là một danh từ dùng để chỉ quận, huyện hoặc một khu vực đặc thù.}

\vocabentry{Plaza}
{/ˈplɑːzə/ (n)}
{A public square, marketplace, or similar open space in a built-up area.}
{The central plaza is a popular meeting place for locals and tourists alike.}
{Đây là một danh từ dùng để chỉ quảng trường hoặc khu vực trống công cộng.}

\vocabentry{Alleyway}
{/ˈæliweɪ/ (n)}
{A narrow passage between or behind buildings.}
{The old town is full of mysterious alleyways and hidden courtyards.}
{Đây là một danh từ dùng để chỉ con hẻm nhỏ.}

\vocabentry{Slum clearance}
{/slʌm ˈklɪərəns/ (n)}
{The removal of slums to make way for new buildings.}
{Slum clearance programs must ensure that the residents are provided with alternative housing.}
{Đây là một cụm danh từ dùng để chỉ việc giải tỏa khu ổ chuột.}

\vocabentry{Urban decay}
{/ˌɜːrbən dɪˈkeɪ/ (n)}
{The process by which a previously functioning city, or part of a city, falls into disrepair.}
{Urban decay is visible in parts of the city where factories have closed down.}
{Đây là một cụm danh từ dùng để chỉ sự suy tàn đô thị.}

\vocabentry{Social unrest}
{/ˌsoʊʃl ʌnˈrest/ (n)}
{A state of dissatisfaction and disturbance among groups of people.}
{Economic inequality in cities can sometimes lead to violent social unrest.}
{Đây là một cụm danh từ dùng để chỉ sự bất ổn xã hội, bạo động.}

\vocabentry{Segregation}
{/ˌseɡrɪˈɡeɪʃn/ (n)}
{The action or state of setting someone or something apart from others. (Đây là một danh từ dùng để chỉ sự phân biệt đối xử hoặc tách biệt dân cư (giàu/nghèo, sắc tộc).).}
{Social segregation is still a significant issue in many large American cities.}
{}

\vocabentry{Integration}
{/ˌɪntɪˈɡreɪʃn/ (n)}
{The action or process of integrating.}
{The integration of new immigrants into the urban community is a slow process.}
{Đây là một danh từ dùng để chỉ sự hòa nhập cộng đồng.}

\vocabentry{Inclusivity}
{/ˌɪnkluːˈsɪvəti/ (n)}
{The practice or policy of including people who might otherwise be excluded or marginalized.}
{Urban planners are focusing more on inclusivity to ensure that cities work for everyone.}
{Đây là một danh từ dùng để chỉ tính hòa nhập, không phân biệt đối xử.}

\vocabentry{Sustainability}
{/səˌsteɪnəˈbɪləti/ (n)}
{The ability to be maintained at a certain rate or level.}
{Environmental sustainability is a key goal for modern city development.}
{Đây là một danh từ dùng để chỉ tính bền vững của đô thị về mặt môi trường và kinh tế.}

\vocabentry{Walkability}
{/ˌwɔːkəˈbɪləti/ (n)}
{A measure of how friendly an area is to walking.}
{High walkability reduces the need for cars and promotes a healthier lifestyle.}
{Đây là một danh từ dùng để chỉ tính thuận tiện cho việc đi bộ.}

\vocabentry{Mixed-income}
{/ˌmɪkst ˈɪnkʌm/ (adj)}
{Consisting of households with various income levels.}
{Mixed-income housing developments can help prevent the formation of ghettos.}
{Đây là một tính từ dùng để chỉ các khu dân cư có nhiều mức thu nhập khác nhau, tránh phân cực giàu nghèo.}

\vocabentry{Public utility}
{/ˌpʌblɪk juːˈtɪləti/ (n)}
{An organization supplying the community with electricity, gas, water, or sewerage.}
{Public utilities must be reliable and affordable for all city residents.}
{Đây là một cụm danh từ dùng để chỉ các dịch vụ tiện ích công cộng.}

\vocabentry{Communal}
{/kəˈmjuːnl/ (adj)}
{Shared by all members of a community; for common use.}
{The apartment building has a communal garden where residents can relax.}
{Đây là một tính từ dùng để chỉ các không gian hoặc tiện ích chung.}

\vocabentry{Proximity}
{/prɑːkˈsɪməti/ (n)}
{Nearness in space, time, or relationship.}
{The proximity of the school to their house was a major factor in their decision to buy it.}
{Đây là một danh từ dùng để chỉ sự gần gũi về khoảng cách địa lý.}

\vocabentry{Upward mobility}
{/ˌʌpwərd moʊˈbɪləti/ (n)}
{Rising from a lower to a higher social class or status.}
{Cities are often seen as engines of upward mobility for ambitious young people.}
{Đây là một cụm danh từ dùng để chỉ sự thăng tiến trong xã hội.}

\vocabentry{Urban sprawl}
{/ˌɜːrbən ˈsprɔːl/ (n)}
{The uncontrolled expansion of urban areas.}
{Urban sprawl increases travel times and consumes valuable farmland.}
{Đây là một cụm danh từ dùng để chỉ sự mở rộng đô thị không kiểm soát.}

\vocabentry{Car-sharing}
{/ˈkɑːr ˌʃerɪŋ/ (n)}
{An arrangement in which several people travel together in one car.}
{Car-sharing is an effective way to reduce the number of vehicles on the road.}
{Đây là một danh từ dùng để chỉ hình thức đi chung xe.}

\vocabentry{Telecommuting}
{/ˌtelikəˈmjuːtɪŋ/ (n)}
{Working from home using the internet and phone.}
{Telecommuting can significantly reduce traffic congestion during peak hours.}
{Đây là một danh từ dùng để chỉ việc làm việc từ xa, không cần đến văn phòng.}

\vocabentry{Smart city}
{/ˈsmɑːrt ˈsɪti/ (n)}
{An urban area that uses different types of electronic methods and sensors to collect data.}
{Smart cities use technology to manage traffic, energy, and waste more efficiently.}
{Đây là một cụm danh từ dùng để chỉ thành phố thông minh.}

\vocabentry{Intelligent transport}
{/ɪnˈtelɪdʒənt ˈtrænspɔːrt/ (n)}
{Advanced applications which aim to provide innovative services relating to transport.}
{Intelligent transport systems can provide real-time information about bus arrival times.}
{Đây là một cụm danh từ dùng để chỉ hệ thống giao thông thông minh.}

\vocabentry{Urban forestry}
{/ˌɜːrbən ˈfɔːrɪstri/ (n)}
{The care and management of single trees and tree populations in urban settings.}
{Urban forestry helps improve air quality and provide shade for pedestrians.}
{Đây là một cụm danh từ dùng để chỉ lâm nghiệp đô thị, trồng cây trong thành phố.}

\vocabentry{Rooftop garden}
{/ˈruːftɑːp ˌɡɑːrdn/ (n)}
{A garden on the roof of a building.}
{Rooftop gardens help insulate buildings and reduce the heat island effect.}
{Đây là một cụm danh từ dùng để chỉ vườn trên sân thượng.}

\vocabentry{Waterfront development}
{/ˈwɔːtərfrʌnt dɪˈveləpmənt/ (n)}
{The redevelopment of land adjacent to a body of water.}
{Waterfront development often attracts high-end restaurants and luxury apartments.}
{Đây là một cụm danh từ dùng để chỉ việc phát triển các khu vực ven sông/biển.}

\vocabentry{Public art}
{/ˌpʌblɪk ˈɑːrt/ (n)}
{Art in any media that has been planned and executed with the intention of being staged in the physical public domain. (Đây là một cụm danh từ dùng để chỉ nghệ thuật công cộng (tượng đài, tranh tường).).}
{Public art can help create a unique identity for a neighborhood.}
{}

\vocabentry{Cultural heritage}
{/ˌkʌltʃərəl ˈherɪtɪdʒ/ (n)}
{The legacy of physical artifacts and intangible attributes of a group or society.}
{Preserving cultural heritage is essential during the process of urban modernization.}
{Đây là một cụm danh từ dùng để chỉ di sản văn hóa.}

\vocabentry{Heritage site}
{/ˈherɪtɪdʒ saɪt/ (n)}
{A place that has been registered by the government as having cultural or historical significance.}
{The old temple is a protected heritage site, so it cannot be demolished.}
{Đây là một cụm danh từ dùng để chỉ các khu di tích/di sản.}

\vocabentry{Preservation}
{/ˌprezərˈveɪʃn/ (n)}
{The action of preserving something. (Đây là một danh từ dùng để chỉ sự bảo tồn (các công trình cũ, văn hóa).).}
{The preservation of historic buildings is a key part of urban identity.}
{}

\vocabentry{Gourmet}
{/ˈɡʊrmeɪ/ (adj)}
{Relating to good food and drink.}
{The city is famous for its gourmet restaurants and vibrant food scene.}
{Đây là một tính từ dùng để chỉ sự sành ăn, thường dùng để tả các khu phố ẩm thực cao cấp ở thành phố.}

\vocabentry{Street furniture}
{/ˈstriːt ˌfɜːrnɪtʃər/ (n)}
{Objects and pieces of equipment installed along streets and roads for various purposes. (Đây là một cụm danh từ dùng để chỉ các thiết bị đường phố (ghế băng, cột đèn, thùng rác).).}
{Comfortable street furniture encourages people to spend more time outdoors.}
{}

\vocabentry{Wayfinding}
{/ˈweɪfaɪndɪŋ/ (n)}
{Information systems that guide people through a physical environment and enhance their understanding and experience of the space.}
{Clear wayfinding signs are essential for tourists in a large, unfamiliar city.}
{Đây là một danh từ dùng để chỉ hệ thống chỉ đường.}

\vocabentry{Public health}
{/ˌpʌblɪk ˈhelθ/ (n)}
{The health of the population as a whole.}
{Urban planning plays a crucial role in promoting public health by encouraging exercise.}
{Đây là một cụm danh từ dùng để chỉ sức khỏe cộng đồng.}

\vocabentry{Economic engine}
{/ˌiːkəˈnɑːmɪk ˈendʒɪn/ (n)}
{Something that provides a lot of jobs and money for an area.}
{Large cities serve as the primary economic engines of their respective countries.}
{Đây là một cụm danh từ dùng để chỉ động lực kinh tế.}

\vocabentry{Gentrify}
{/ˈdʒentrɪfaɪ/ (v)}
{Renovate and improve (especially a house or district) so that it conforms to middle-class taste.}
{The area began to gentrify as young artists moved into the old warehouses.}
{Đây là một động từ dùng để chỉ việc chỉnh trang khu dân cư cho tầng lớp trung lưu.}

\vocabentry{Slumlord}
{/ˈslʌmlɔːrd/ (n)}
{A landlord who receives unusually high rent from tenants in poorly maintained slum properties.}
{The city is taking legal action against slumlords who ignore safety regulations.}
{Đây là một danh từ chỉ những chủ nhà cho thuê ở khu ổ chuột với giá cao nhưng không sửa chữa.}

\vocabentry{Homelessness}
{/ˈhoʊmləsnəs/ (n)}
{The state of having no home.}
{Homelessness is a growing problem in many cities due to rising rents.}
{Đây là một danh từ dùng để chỉ tình trạng vô gia cư.}

\vocabentry{Deprivation}
{/ˌdeprɪˈveɪʃn/ (n)}
{The damaging lack of material benefits considered to be basic necessities in a society.}
{Many inner-city areas suffer from high levels of social and economic deprivation.}
{Đây là một danh từ dùng để chỉ sự thiếu thốn, nghèo khổ.}

\vocabentry{Marginalized}
{/ˈmɑːrdʒɪnəlaɪzd/ (adj)}
{(Of a person, group, or concept) treated as insignificant or peripheral.}
{The government needs to do more to support marginalized communities in urban centers.}
{Đây là một tính từ dùng để chỉ các nhóm người bị gạt ra ngoài lề xã hội.}

\vocabentry{Disadvantaged}
{/ˌdɪsədˈvæntɪdʒd/ (adj)}
{Lacking the money, education, or status to be very successful or happy.}
{Programs are being developed to help disadvantaged youth find jobs in the city.}
{Đây là một tính từ dùng để chỉ những người chịu thiệt thòi về mặt kinh tế xã hội.}

\vocabentry{Underprivileged}
{/ˌʌndərˈprɪvəlɪdʒd/ (adj)}
{(Of a person) not enjoying the same standard of living or rights as the majority of people in a society.}
{The charity provides free meals for underprivileged children in the slums.}
{Đây là một tính từ dùng để chỉ những người không có đủ điều kiện/đặc quyền như số đông.}

\vocabentry{Segregated}
{/ˈseɡrɪɡeɪtɪd/ (adj)}
{Set apart from the rest or from each other; isolated or divided.}
{Some cities remain highly segregated by income and ethnicity.}
{Đây là một tính từ dùng để chỉ sự tách biệt dân cư.}

\vocabentry{Affluence}
{/ˈæfluəns/ (n)}
{The state of having a great deal of money; wealth.}
{The affluence of the suburbs contrasts sharply with the poverty of the inner city.}
{Đây là một danh từ dùng để chỉ sự giàu có, sung túc.}

\vocabentry{Affluent}
{/ˈæfluənt/ (adj)}
{(Especially of a group or area) having a great deal of money; wealthy.}
{They live in an affluent neighborhood where the houses cost millions.}
{Đây là một tính từ dùng để chỉ sự giàu có.}

\vocabentry{Amenities}
{/əˈmenəti/ (n)}
{The useful features or facilities of a building or place. (Đây là một danh từ (số nhiều) dùng để chỉ các tiện ích công cộng.).}
{The city offers many amenities, such as excellent museums and theaters.}
{}

\vocabentry{Automobile}
{/ˈɔːtəməbiːl/ (n)}
{A car.}
{The invention of the automobile led to the growth of the suburbs.}
{Đây là một danh từ dùng để chỉ ô tô, phương tiện gây ra nhiều vấn đề đô thị.}

\vocabentry{Car-centric}
{/kɑːr ˈsentrɪk/ (adj)}
{Focused on or prioritizing the needs of cars.}
{Many American cities are criticized for being too car-centric and unfriendly to pedestrians.}
{Đây là một tính từ dùng để chỉ các quy hoạch ưu tiên cho ô tô thay vì con người.}

\vocabentry{Carpooling}
{/ˈkɑːrpuːlɪŋ/ (n)}
{The sharing of car journeys so that more than one person travels in a car.}
{Carpooling can reduce commuting costs and help the environment.}
{Đây là một danh từ dùng để chỉ hình thức đi chung xe.}

\vocabentry{Congested}
{/kənˈdʒestɪd/ (adj)}
{(Of a road or place) so crowded with traffic or people as to hinder or prevent freedom of movement.}
{The main highway is always congested during the morning rush hour.}
{Đây là một tính từ dùng để chỉ tình trạng tắc nghẽn.}

\vocabentry{Expressway}
{/ɪkˈspresweɪ/ (n)}
{A highway designed for fast traffic.}
{The new expressway has cut the journey time between the two cities in half.}
{Đây là một danh từ dùng để chỉ đường cao tốc.}

\vocabentry{Interchange}
{/ˈɪntərtʃeɪndʒ/ (n)}
{A complex junction where several roads or motorways meet.}
{Traffic often slows down at the busy interchange near the city limits.}
{Đây là một danh từ dùng để chỉ nút giao thông lập thể, nơi nhiều con đường gặp nhau.}

\vocabentry{Transit-oriented}
{/ˈtrænzɪt ˈɔːriəntɪd/ (adj)}
{Focusing on creating a vibrant, walkable, mixed-use community near a transit station.}
{Transit-oriented development helps reduce dependence on cars and promotes sustainable living.}
{Đây là một tính từ dùng để chỉ quy hoạch tập trung quanh các nút giao thông công cộng.}

\vocabentry{Traffic jam}
{/ˈtræfɪk dʒæm/ (n)}
{A line of vehicles that cannot move or can only move very slowly.}
{We were stuck in a traffic jam for over two hours.}
{Đây là một cụm danh từ dùng để chỉ sự tắc nghẽn giao thông.}

\vocabentry{Rush hour}
{/ˈrʌʃ aʊər/ (n)}
{The times at the beginning and end of the working day when many people are traveling to or from work.}
{I try to avoid driving during the rush hour to save time and stress.}
{Đây là một cụm danh từ dùng để chỉ giờ cao điểm.}

\vocabentry{Peak hour}
{/ˈpiːk aʊər/ (n)}
{The time of day when a lot of people use a particular service.}
{Train tickets are more expensive during peak hours.}
{Đây là một cụm danh từ dùng để chỉ giờ cao điểm sử dụng dịch vụ hoặc điện.}

\vocabentry{Urbanites}
{/ˈɜːrbənaɪts/ (n)}
{People who live in a city or town. (Đây là một danh từ (số nhiều) dùng để chỉ dân thành thị.).}
{Urbanites often appreciate the convenience and variety of city life.}
{}

\vocabentry{Apartment block}
{/əˈpɑːrtmənt blɑːk/ (n)}
{A large building containing many apartments.}
{A new apartment block is being built in the city center.}
{Đây là một cụm danh từ dùng để chỉ khu chung cư.}

\vocabentry{Condominium}
{/ˌkɑːndəˈmɪniəm/ (n)}
{A building or complex of buildings containing a number of individually owned apartments or houses.}
{He bought a luxury condominium with a view of the harbor.}
{Đây là một danh từ dùng để chỉ căn hộ chung cư cao cấp.}

\vocabentry{Residential area}
{/ˌrezɪˈdenʃl ˈeriə/ (n)}
{An area where people live.}
{This is a quiet residential area, so please keep the noise down.}
{Đây là một cụm danh từ dùng để chỉ khu dân cư.}

\vocabentry{Borough}
{/ˈbɜːroʊ/ (n)}
{A town or district which is an administrative unit.}
{New York City is made up of five boroughs: Manhattan, Brooklyn, Queens, the Bronx, and Staten Island.}
{Đây là một danh từ dùng để chỉ một quận hoặc đơn vị hành chính của thành phố.}

\vocabentry{Precinct}
{/ˈpriːsɪŋkt/ (n)}
{An area within a city or town, typically one that is closed to traffic.}
{The shopping precinct is always crowded on Saturdays.}
{Đây là một danh từ dùng để chỉ phân khu hoặc khu vực bộ hành.}

\vocabentry{Ward}
{/wɔːrd/ (n)}
{A district into which a city or town is divided for the purpose of administration and elections.}
{Each ward has its own representative on the city council.}
{Đây là một danh từ dùng để chỉ phường hoặc đơn vị bầu cử nhỏ.}

\vocabentry{Municipality}
{/mjuːˌnɪsɪˈpæləti/ (n)}
{A city or town that has corporate status and local government.}
{The municipality is responsible for waste collection and street cleaning.}
{Đây là một danh từ dùng để chỉ đô thị tự trị hoặc chính quyền đô thị.}

\vocabentry{Town hall}
{/ˌtaʊn ˈhɔːl/ (n)}
{A building used for the administration of local government.}
{The protest took place outside the town hall.}
{Đây là một cụm danh từ dùng để chỉ tòa thị chính.}

\vocabentry{Civic}
{/ˈsɪvɪk/ (adj)}
{Relating to a city or town, especially its administration; municipal.}
{Good citizens take an active interest in civic affairs.}
{Đây là một tính từ dùng để chỉ các vấn đề thuộc về nghĩa vụ hoặc hoạt động dân sự đô thị.}

\vocabentry{Bureaucracy}
{/bjʊˈrɑːkrəsi/ (n)}
{A system of government in which most of the important decisions are made by state officials rather than by elected representatives.}
{You have to deal with a lot of bureaucracy to get a building permit.}
{Đây là một danh từ dùng để chỉ bộ máy hành chính cồng kềnh, thường gặp ở các thành phố lớn.}

\vocabentry{Metropolitan}
{/ˌmetrəˈpɑːlɪtən/ (adj)}
{Relating to or denoting a metropolis.}
{The metropolitan area has a population of over five million people.}
{Đây là một tính từ dùng để chỉ những gì thuộc về vùng đô thị lớn.}

\vocabentry{Suburban}
{/səˈbɜːrbən/ (adj)}
{Relating to or characteristic of a suburb.}
{He finds suburban life a bit boring and misses the excitement of the city.}
{Đây là một tính từ dùng để chỉ các đặc điểm của vùng ngoại ô.}

\vocabentry{Exurban}
{/eksˈɜːrbən/ (adj)}
{Relating to an exurb (a region beyond the suburbs).}
{Exurban development is a relatively new phenomenon in this country.}
{Đây là một tính từ dùng để chỉ vùng xa hơn ngoại ô, nơi dân thành phố dời về để tìm sự yên tĩnh tuyệt đối.}

\vocabentry{Commercialized}
{/kəˈmɜːrʃəlaɪzd/ (adj)}
{Managed or exploited in a way designed to make a profit.}
{The historic district has become too commercialized, with souvenir shops everywhere.}
{Đây là một tính từ dùng để chỉ sự thương mại hóa quá mức.}

\vocabentry{Industrialized}
{/ɪnˈdʌstriəlaɪzd/ (adj)}
{(Of a country or region) having developed industries on a wide scale.}
{Most industrialized nations are now facing the challenge of an aging population.}
{Đây là một tính từ dùng để chỉ sự công nghiệp hóa.}

\vocabentry{Landscaped}
{/ˈlændskeɪpt/ (adj)}
{(Of a garden or area of land) laid out with plants, trees, etc. to make it more attractive.}
{The new office park features beautifully landscaped gardens and fountains.}
{Đây là một tính từ dùng để chỉ các khu vực được thiết kế cảnh quan.}

\vocabentry{Greenbelt}
{/ˈɡriːnbelt/ (n)}
{An area of open land around a city, on which building is restricted.}
{The greenbelt protects the city from uncontrolled urban sprawl.}
{Đây là một danh từ dùng để chỉ vành đai xanh, khu vực hạn chế xây dựng quanh thành phố.}

\vocabentry{Public space}
{/ˌpʌblɪk ˈspeɪs/ (n)}
{A social space that is generally open and accessible to people.}
{Parks, plazas, and libraries are all important types of public space.}
{Đây là một cụm danh từ dùng để chỉ không gian công cộng.}

\vocabentry{Brownfield}
{/ˈbraʊnfiːld/ (n)}
{A former industrial or commercial site where future use is affected by environmental contamination. (Đây là một danh từ dùng để chỉ các khu đất cũ (nhà máy, kho bãi) bị bỏ hoang chờ tái quy hoạch.).}
{The city is converting old brownfield sites into new housing and parks.}
{}

\vocabentry{Greenfield}
{/ˈɡriːnfiːld/ (n)}
{Denoting or relating to previously undeveloped land.}
{Developers often prefer greenfield sites because they are easier and cheaper to build on.}
{Đây là một danh từ dùng để chỉ các khu đất trống chưa từng xây dựng, thường là đất nông nghiệp.}

\vocabentry{Built-up}
{/ˌbɪlt ˈʌp/ (adj)}
{(Of an area) having a lot of buildings.}
{It's hard to find a quiet place in such a built-up city.}
{Đây là một tính từ dùng để chỉ khu vực đã được xây dựng san sát.}

\vocabentry{Densely populated}
{/ˌdensli ˈpɑːpjuletɪd/ (adj)}
{Having a lot of people living in a small area.}
{Java is one of the most densely populated islands in the world.}
{Đây là một cụm tính từ dùng để chỉ mật độ dân cư đông đúc.}

\vocabentry{Sparsely populated}
{/ˌspɑːrsli ˈpɑːpjuletɪd/ (adj)}
{Having few people living in a large area.}
{The northern part of Canada is sparsely populated.}
{Đây là một cụm tính từ dùng để chỉ mật độ dân cư thưa thớt.}

\vocabentry{Uninhabitable}
{/ˌʌnɪnˈhæbɪtəbl/ (adj)}
{(Of a place) unsuitable for living in.}
{Many buildings were rendered uninhabitable by the earthquake.}
{Đây là một tính từ dùng để chỉ những nơi không thể ở được.}

\vocabentry{Residential complex}
{/ˌrezɪˈdenʃl ˈkɑːmpleks/ (n)}
{A group of residential buildings.}
{The new residential complex has its own shopping center and cinema.}
{Đây là một cụm danh từ dùng để chỉ tổ hợp nhà ở/chung cư.}

\vocabentry{Shopping mall}
{/ˈʃɑːpɪŋ mɔːl/ (n)}
{A large building or group of buildings containing many stores.}
{On weekends, the shopping mall is packed with families.}
{Đây là một cụm danh từ dùng để chỉ trung tâm thương mại.}

\vocabentry{Entertainment hub}
{/ˌentərˈteɪnmənt hʌb/ (n)}
{A central area for movies, theaters, and other activities.}
{The area near the river has become a major entertainment hub for the city.}
{Đây là một cụm danh từ dùng để chỉ trung tâm vui chơi giải trí.}

\vocabentry{Financial district}
{/faɪˈnænʃl ˈdɪstrɪkt/ (n)}
{An area where many banks and financial companies are located.}
{Wall Street is the most famous financial district in the world.}
{Đây là một cụm danh từ dùng để chỉ khu trung tâm tài chính.}

\vocabentry{Nightspot}
{/ˈnaɪtspɑːt/ (n)}
{A place for entertainment at night, such as a club or bar.}
{The city is full of trendy nightspots that are popular with young people.}
{Đây là một danh từ dùng để chỉ địa điểm vui chơi về đêm.}

\vocabentry{Vibrant}
{/ˈvaɪbrənt/ (adj)}
{Full of energy and life.}
{I love the vibrant atmosphere of the night markets.}
{Đây là một tính từ dùng để chỉ sự náo nhiệt, tràn đầy năng lượng.}

\vocabentry{Atmospheric}
{/ˌætməsˈferɪk/ (adj)}
{Creating a distinctive mood or environment. (Đây là một tính từ dùng để chỉ một nơi có không gian/không khí đặc trưng (thường là cổ kính hoặc lãng mạn).).}
{The old town is very atmospheric, especially at night when the streets are lit by gas lamps.}
{}

\vocabentry{Picturesque}
{/ˌpɪktʃəˈresk/ (adj)}
{Visually attractive, especially in a quaint or charming way. (Đây là một tính từ dùng để chỉ vẻ đẹp đẹp như tranh vẽ (thường cho các thị trấn cổ).).}
{We spent the day wandering through the picturesque streets of the village.}
{}

\vocabentry{Quaint}
{/kweɪnt/ (adj)}
{Attractively unusual or old-fashioned.}
{The town is full of quaint little shops and cafes.}
{Đây là một tính từ dùng để chỉ vẻ đẹp cổ kính, độc lạ một cách quyến rũ.}

\vocabentry{Dilapidated}
{/dɪˈlæpɪdeɪtɪd/ (adj)}
{(Of a building or object) in a state of disrepair or ruin as a result of age or neglect.}
{The city is working to renovate or demolish dilapidated buildings in the inner city.}
{Đây là một tính từ dùng để chỉ những tòa nhà cũ kỹ, đổ nát, xuống cấp.}

\vocabentry{Derelict}
{/ˈderəlɪkt/ (adj)}
{In a very poor condition as a result of disuse and neglect.}
{The derelict factory has become a haven for homeless people and drug addicts.}
{Đây là một tính từ dùng để chỉ các khu vực bị bỏ hoang, không ai chăm sóc.}

\vocabentry{Run-down}
{/ˌrʌn ˈdaʊn/ (adj)}
{(Of a building or area) in a poor or neglected state.}
{They live in a run-down part of the city where the streets are dirty and the buildings need repair.}
{Đây là một tính từ dùng để chỉ sự xuống cấp, tồi tàn của một khu phố.}

\vocabentry{Squalid}
{/ˈskwɑːlɪd/ (adj)}
{(Of a place) extremely dirty and unpleasant, especially as a result of poverty or neglect.}
{Thousands of people are forced to live in squalid conditions in the slums.}
{Đây là một tính từ dùng để chỉ sự dơ bẩn, bệ rạc do nghèo khổ.}

\vocabentry{Filthy}
{/ˈfɪlθi/ (adj)}
{Disgustingly dirty.}
{The city's rivers have become filthy due to industrial waste and sewage.}
{Đây là một tính từ cực cấp dùng để chỉ sự bẩn thỉu.}

\vocabentry{Grime}
{/ɡraɪm/ (n)}
{Dirt ingrained on the surface of something.}
{The windows were covered in years of urban grime.}
{Đây là một danh từ dùng để chỉ lớp bụi bẩn bám lâu ngày, đặc biệt ở các thành phố công nghiệp.}

\vocabentry{Littered}
{/ˈlɪtərd/ (adj)}
{Untidy with rubbish or a large number of objects left lying about.}
{The park was littered with plastic bottles and food wrappers after the concert.}
{Đây là một tính từ dùng để chỉ nơi tràn ngập rác rưởi.}

\vocabentry{Graffiti}
{/ɡrəˈfiːti/ (n)}
{Writing or drawings scribbled, scratched, or sprayed illicitly on a wall or other surface in a public place.}
{The subway walls are covered in colorful but illegal graffiti.}
{Đây là một danh từ dùng để chỉ hình vẽ bậy trên tường.}

\vocabentry{Vandalism}
{/ˈvændəlɪzəm/ (n)}
{Action involving deliberate destruction of or damage to public or private property.}
{The city is spending millions of dollars to repair the damage caused by vandalism.}
{Đây là một danh từ dùng để chỉ hành vi phá hoại tài sản công cộng.}

\vocabentry{Antisocial}
{/ˌæntiˈsoʊʃl/ (adj)}
{Contrary to the laws and customs of society; devoid of or antagonistic to sociable instincts or practices.}
{Graffiti and loud noise are often considered forms of antisocial behavior.}
{Đây là một tính từ dùng để chỉ các hành vi gây rối trật tự công cộng.}

\vocabentry{Law-abiding}
{/ˈlɔː əbaɪdɪŋ/ (adj)}
{Obedient to the laws of society.}
{The majority of city residents are law-abiding citizens who want a safe environment.}
{Đây là một tính từ dùng để chỉ những người dân tuân thủ pháp luật.}

\vocabentry{Safe}
{/seɪf/ (adj)}
{Protected from or not exposed to danger or risk.}
{It is generally safe to walk around the city center at night, but you should still be careful.}
{Đây là một tính từ dùng để chỉ tính an toàn.}

\vocabentry{Dangerous}
{/ˈdeɪndʒərəs/ (adj)}
{Able or likely to cause harm or injury.}
{Some parts of the city are considered dangerous due to high crime rates.}
{Đây là một tính từ dùng để chỉ sự nguy hiểm.}

\vocabentry{CCTV}
{/ˌsiː siː tiː ˈviː/ (n)}
{Closed-circuit television.}
{The city has installed CCTV cameras in public squares to deter crime.}
{Đây là một danh từ dùng để chỉ hệ thống camera giám sát an ninh.}

\vocabentry{Policing}
{/pəˈliːsɪŋ/ (n)}
{The maintenance of law and order by a police force.}
{Community policing can help improve the relationship between the police and local residents.}
{Đây là một danh từ dùng để chỉ hoạt động tuần tra, giữ gìn trật tự của cảnh sát.}

\vocabentry{Patrol}
{/pəˈtroʊl/ (v)}
{Keep watch over (an area) by regularly walking or traveling around or through it.}
{Police officers patrol the city streets on foot and in cars to prevent crime.}
{Đây là một động từ dùng để chỉ việc tuần tra.}

\vocabentry{Insecurity}
{/ˌɪnsɪˈkjʊrəti/ (n)}
{The state of being open to danger or threat; lack of protection.}
{Economic insecurity can lead to higher crime rates in urban areas.}
{Đây là một danh từ dùng để chỉ sự bất an, không an toàn.}

\vocabentry{Social cohesion}
{/ˌsoʊʃl koʊˈhiːʒn/ (n)}
{The bonds or "glue" that bring people together in a society.}
{Public events and festivals are important for promoting social cohesion in diverse cities.}
{Đây là một cụm danh từ dùng để chỉ sự gắn kết xã hội.}

\vocabentry{Sense of belonging}
{/sens əv bɪˈlɔːŋɪŋ/ (n)}
{A feeling that you are a part of a group or a place.}
{Having a sense of belonging is essential for the mental health of city residents.}
{Đây là một cụm danh từ dùng để chỉ cảm giác thuộc về nơi mình đang sống.}

\vocabentry{Alienation}
{/ˌeɪliəˈneɪʃn/ (n)}
{The state or experience of being isolated from a group or an activity to which one should belong or in which one should be involved.}
{The anonymity of the city can sometimes lead to feelings of alienation and loneliness.}
{Đây là một danh từ dùng để chỉ sự lạc lõng, ghẻ lạnh giữa đám đông đô thị.}

\vocabentry{Anonymity}
{/ˌænəˈnɪməti/ (n)}
{The condition of being anonymous.}
{Many people appreciate the anonymity of city life, as it gives them more freedom.}
{Đây là một danh từ dùng để chỉ sự vô danh, không ai biết ai ở thành phố lớn.}

\vocabentry{Identity}
{/aɪˈdentəti/ (n)}
{The fact of being who or what a person or thing is.}
{Each city neighborhood has its own unique identity and atmosphere.}
{Đây là một danh từ dùng để chỉ bản sắc hoặc đặc trưng riêng của một thành phố.}

\vocabentry{Multiculturalism}
{/ˌmʌltiˈkʌltʃərəlɪzəm/ (n)}
{The presence of, or support for the presence of, several distinct cultural or ethnic groups within a society.}
{Multiculturalism is a source of strength and creativity for global cities.}
{Đây là một danh từ dùng để chỉ chủ nghĩa đa văn hóa.}

\vocabentry{Integration}
{/ˌɪntɪˈɡreɪʃn/ (n)}
{The process of becoming part of a group of people.}
{The integration of immigrants into the local community is a key challenge for urban governments.}
{Đây là một danh từ dùng để chỉ sự hòa nhập.}

\vocabentry{Tolerance}
{/ˈtɑːlərəns/ (n)}
{The ability or willingness to tolerate something, in particular the existence of opinions or behavior that one does not necessarily agree with.}
{Large cities usually have a higher level of tolerance for different lifestyles.}
{Đây là một danh từ dùng để chỉ sự bao dung, chấp nhận sự khác biệt ở thành phố.}

\vocabentry{Prejudice}
{/ˈpredʒədɪs/ (n)}
{Preconceived opinion that is not based on reason or actual experience.}
{We must work together to overcome prejudice and discrimination in our urban communities.}
{Đây là một danh từ dùng để chỉ định kiến.}

\vocabentry{Discrimination}
{/dɪˌskrɪmɪˈneɪʃn/ (n)}
{The unjust or prejudicial treatment of different categories of people or things, especially on the grounds of race, age, or sex.}
{Discrimination in the housing market is illegal but still occurs in some areas.}
{Đây là một danh từ dùng để chỉ sự phân biệt đối xử.}

\vocabentry{Equality}
{/iˈkwɑːləti/ (n)}
{The state of being equal, especially in status, rights, and opportunities.}
{Promoting gender equality is essential for the development of modern cities.}
{Đây là một danh từ dùng để chỉ sự bình đẳng.}

\vocabentry{Inequality}
{/ˌɪnɪˈkwɑːləti/ (n)}
{Difference in size, degree, circumstances, etc.; lack of equality. (Đây là một danh từ dùng để chỉ sự bất bình đẳng (thu nhập, cơ hội).).}
{Economic inequality is a major driver of social problems in urban centers.}
{}

\vocabentry{Fairness}
{/ˈfernəs/ (n)}
{Impartial and just treatment or behavior without favoritism or discrimination.}
{The tax system should be designed with fairness in mind to support low-income city residents.}
{Đây là một danh từ dùng để chỉ sự công bằng.}

\vocabentry{Justice}
{/ˈdʒʌstɪs/ (n)}
{Just behavior or treatment.}
{Social justice is a key goal for many urban community organizations.}
{Đây là một danh từ dùng để chỉ công lý.}

\vocabentry{Welfare}
{/ˈwelfer/ (n)}
{The health, happiness, and fortunes of a person or group.}
{The city government provides various welfare services for its most vulnerable residents.}
{Đây là một danh từ dùng để chỉ phúc lợi xã hội.}

\vocabentry{Quality of life}
{/ˌkwɑːləti əv ˈlaɪf/ (n)}
{The standard of health, comfort, and happiness experienced by an individual or group.}
{Improving the quality of life for all residents is the primary goal of urban planning.}
{Đây là một cụm danh từ dùng để chỉ chất lượng cuộc sống.}

\vocabentry{Well-being}
{/ˌwel ˈbiːɪŋ/ (n)}
{The state of being comfortable, healthy, or happy.}
{Access to parks and green spaces is essential for the physical and mental well-being of city dwellers.}
{Đây là một danh từ dùng để chỉ trạng thái hạnh phúc và khỏe mạnh.}

\vocabentry{Liveability}
{/ˌlɪvəˈbɪləti/ (n)}
{The sum of the factors that add up to a community's quality of life—including the built and natural environments, economic prosperity, social stability and equity, educational opportunity, and cultural, entertainment and recreation possibilities.}
{The new public transit system has greatly improved the liveability of the city.}
{Đây là một danh từ dùng để chỉ mức độ đáng sống của một thành phố.}

